\chapter{Optimization of Alignments}

\section{Motivation}

\begin{remark}
In practical applications, alignments are not only reconstructed from
data, but also designed, modified, and improved.
\end{remark}

\begin{remark}
This leads naturally to optimization problems in which alignment
parameters are adjusted to satisfy geometric, physical, and regulatory
constraints.
\end{remark}

\section{General Optimization Problem}

\begin{definition}[Alignment optimization problem]
An alignment optimization problem consists of:
\begin{itemize}
\item a parameterized alignment model,
\item an objective function,
\item a set of constraints.
\end{itemize}
\end{definition}

\begin{definition}[Objective function]
An objective function is a mapping
\[
J : \mathcal{A} \to \R
\]
that assigns a cost or quality measure to an alignment $\mathcal{A}$.
\end{definition}

\begin{remark}
Typical objectives include:
\begin{itemize}
\item minimizing curvature variation,
\item minimizing deviation from measurement data,
\item minimizing construction cost,
\item maximizing comfort.
\end{itemize}
\end{remark}

\section{Parameterization}

\begin{remark}
The choice of parameters is central to the formulation of the
optimization problem.
\end{remark}

\begin{definition}[Parameter vector]
An alignment is described by a parameter vector
\[
\mathbf{x} \in \R^n
\]
encoding quantities such as:
\begin{itemize}
\item element lengths,
\item curvature values,
\item transition parameters,
\item geometric constraints.
\end{itemize}
\end{definition}

\begin{remark}
The sparse alignment model provides a natural parameterization by
element-wise variables.
\end{remark}

\TODO{Explicit parameter sets for fixed and transition elements.}

\section{Constraints}

\begin{definition}[Constraint]
A constraint is a condition of the form
\[
g_i(\mathbf{x}) \leq 0,
\qquad
h_j(\mathbf{x}) = 0.
\]
\end{definition}

\begin{remark}
Constraints arise from:
\begin{itemize}
\item geometric continuity,
\item engineering design rules,
\item physical limitations,
\item regulatory requirements.
\end{itemize}
\end{remark}

\subsection{Geometric Constraints}

\begin{itemize}
\item continuity of position and direction,
\item continuity of curvature.
\end{itemize}

\subsection{Engineering Constraints}

\begin{itemize}
\item minimum radius,
\item maximum cant,
\item limits on curvature derivatives.
\end{itemize}

\OPEN{Formal representation of engineering rules in constraint form.}

\section{Relation to Numerical Reconstruction}

\begin{remark}
Evaluation of the objective function requires reconstruction of the
alignment geometry from the parameter vector.
\end{remark}

\begin{remark}
Thus, optimization is intrinsically linked to numerical reconstruction,
as discussed in the previous chapter.
\end{remark}

\section{Sequential Quadratic Programming}

\begin{remark}
Sequential Quadratic Programming (SQP) is a powerful method for solving
nonlinear constrained optimization problems
\cite{Nocedal2006NumericalOptimization}.
\end{remark}

\begin{definition}[SQP iteration]
An SQP method approximates the original problem by a sequence of
quadratic subproblems.
\end{definition}

\begin{remark}
Each iteration solves a quadratic approximation of the objective
function subject to linearized constraints.
\end{remark}

\TODO{Explicit formulation of SQP for alignment parameters.}

\RESEARCH{Adaptation of SQP to sparse alignment structures.}

\section{Alignment Optimization as Control Problem}

\begin{remark}
The curvature function $\kappaf(s)$ may be interpreted as a control
variable in a continuous system.
\end{remark}

\begin{definition}[Control formulation]
The alignment problem can be formulated as an optimal control problem
with state variables $(\gamma,\theta)$ and control $\kappaf$.
\end{definition}

\RESEARCH{Connection between alignment optimization and optimal control theory.}

\section{Network-Level Optimization}

\begin{remark}
Real-world railway systems involve not only single alignments, but
networks with branching and merging structures.
\end{remark}

\RESEARCH{Extension of optimization from single alignments to networks.}

\OPEN{Handling of switches, crossings, and parallel tracks.}

\section{Robustness and Uncertainty}

\begin{remark}
Optimization must account for uncertainty in input data.
\end{remark}

\RESEARCH{Robust optimization under uncertain geometry.}

\OPEN{Trade-off between optimality and robustness.}

\section{Implementation Perspective}

\begin{remark}
Practical optimization systems must balance:
\begin{itemize}
\item numerical stability,
\item computational efficiency,
\item model expressiveness.
\end{itemize}
\end{remark}

\TODO{Integration of optimization into the modelling workflow.}

\section{Historical Context}

\begin{remark}
Classical systems such as AXTRAN have demonstrated the feasibility of
optimization-based alignment design.
\end{remark}

\begin{remark}
Modern approaches aim to extend these ideas using more flexible models
and improved numerical methods.
\end{remark}

\TODO{Reconstruction and generalization of AXTRAN-like methods.}

\section{Formal Optimization in Curvature Space}

\subsection{Problem Setting}

\begin{remark}
Let \(L>0\) be fixed and let
\[
\kappa : [0,L] \to \R
\]
denote the curvature function of a transition.
\end{remark}

\begin{remark}
We assume the boundary conditions
\[
\kappa(0)=0,
\qquad
\kappa(L)=\kappa_1,
\]
and optionally
\[
\kappa'(0)=0,
\qquad
\kappa'(L)=0.
\]
\end{remark}

\subsection{Admissible Set}

\begin{definition}[Admissible transition set]
\[
\mathcal{A}
=
\left\{
\kappa \in H^1(0,L)
\mid
\kappa(0)=0,\;
\kappa(L)=\kappa_1
\right\}.
\]
\end{definition}

\subsection{Variational Problem}

\begin{definition}[Quadratic curvature-gradient functional]
\[
J[\kappa]
=
\int_0^L \bigl(\kappa'(s)\bigr)^2 \, \dd s .
\]
\end{definition}

\begin{definition}[Optimization problem]
Find $\kappa^\ast \in \mathcal{A}$ such that
\[
J[\kappa^\ast]
=
\min_{\kappa \in \mathcal{A}} J[\kappa].
\]
\end{definition}

\subsection{Euler--Lagrange Equation}

\begin{proposition}
If \(\kappa^\ast\) is a minimizer, then
\[
\kappa^{\ast\prime\prime}(s)=0.
\]
\end{proposition}

\subsection{Interpretation}

\begin{proposition}
The minimizer is
\[
\kappa^\ast(s)=\frac{\kappa_1}{L}s .
\]
\end{proposition}

\begin{remark}
Thus, the clothoid appears as the solution of a variational problem.
\end{remark}

\subsection{Higher-Order Model}

\begin{definition}
\[
J_2[\kappa]
=
\int_0^L \bigl(\kappa''(s)\bigr)^2 \,\dd s.
\]
\end{definition}

\begin{proposition}
Minimizers satisfy
\[
\kappa^{(4)}(s)=0.
\]
\end{proposition}

\begin{remark}
This explains the occurrence of cubic transition laws.
\end{remark}

\subsection{Generalized View}

\begin{remark}
General functionals of the form
\[
J[\kappa]
=
\int_0^L
F\bigl(\kappa,\kappa',\kappa'';v,c\bigr)\,\dd s
\]
allow incorporation of dynamic and engineering criteria.
\end{remark}

\subsection{Connection to the Main Thesis}

\begin{summary}
Transition design can be formulated as an optimization problem in
curvature space.

Classical transition curves arise as solutions of such problems,
supporting the interpretation of transitions as curvature control
functions.
\end{summary}

\section{Coupled Optimization of Curvature and Cant}

\subsection{Motivation}

\begin{remark}
In railway applications, curvature and cant cannot be optimized
independently.
Both jointly determine the effective lateral acceleration and its rate
of change.
\end{remark}

\begin{remark}
This motivates a coupled optimization problem in which both
\[
\kappa(s)
\quad \text{and} \quad
\phi(s)
\]
are treated as design variables.
\end{remark}

\subsection{Effective Acceleration Functional}

\begin{definition}[Effective lateral acceleration]
The effective lateral acceleration is
\[
a_{\mathrm{eff}}(s)
=
v^2 \kappa(s) - g \sin\phi(s).
\]
\end{definition}

\begin{remark}
A natural optimization objective is to reduce the magnitude of
unbalanced lateral acceleration.
\end{remark}

\begin{definition}[Acceleration-based functional]
Define
\[
J_a[\kappa,\phi]
=
\int_0^L a_{\mathrm{eff}}(s)^2 \,\dd s
=
\int_0^L
\bigl(v^2 \kappa(s) - g \sin\phi(s)\bigr)^2
\,\dd s.
\]
\end{definition}

\subsection{Jerk-Based Functional}

\begin{definition}[Jerk]
The jerk is given by
\[
j(s)
=
v^3 \kappa'(s) - g v \cos\phi(s)\,\phi'(s).
\]
\end{definition}

\begin{definition}[Jerk functional]
Define
\[
J_j[\kappa,\phi]
=
\int_0^L j(s)^2 \,\dd s.
\]
That is,
\[
J_j[\kappa,\phi]
=
\int_0^L
\bigl(
v^3 \kappa'(s) - g v \cos\phi(s)\,\phi'(s)
\bigr)^2
\,\dd s.
\]
\end{definition}

\begin{remark}
This functional penalizes rapid variation of effective acceleration and
therefore provides a natural dynamic smoothness criterion.
\end{remark}

\subsection{Combined Objective}

\begin{definition}[Coupled objective functional]
A general coupled objective may be written as
\[
J[\kappa,\phi]
=
\alpha J_a[\kappa,\phi]
+
\beta J_j[\kappa,\phi]
+
\gamma J_\kappa[\kappa]
+
\delta J_\phi[\phi],
\]
where:
\[
J_\kappa[\kappa]
=
\int_0^L (\kappa'(s))^2 \,\dd s,
\qquad
J_\phi[\phi]
=
\int_0^L (\phi'(s))^2 \,\dd s.
\]
\end{definition}

\begin{remark}
The weights
\[
\alpha,\beta,\gamma,\delta \ge 0
\]
determine the relative importance of equilibrium, comfort, curvature
smoothness, and cant smoothness.
\end{remark}

\subsection{Admissible Set}

\begin{definition}[Admissible railway state set]
Let
\[
\mathcal{B}
=
\left\{
(\kappa,\phi)
\;\middle|\;
\kappa \in H^1(0,L),\;
\phi \in H^1(0,L),
\;
\kappa(0)=0,\;
\kappa(L)=\kappa_1
\right\}.
\]
\end{definition}

\begin{remark}
Additional boundary conditions may be imposed, for example:
\[
\phi(0)=0,
\qquad
\phi(L)=\phi_1,
\]
or derivative constraints such as
\[
\kappa'(0)=\kappa'(L)=0,
\qquad
\phi'(0)=\phi'(L)=0.
\]
\end{remark}

\subsection{Coupled Optimization Problem}

\begin{definition}[Coupled railway optimization problem]
Find
\[
(\kappa^\ast,\phi^\ast)\in\mathcal{B}
\]
such that
\[
J[\kappa^\ast,\phi^\ast]
=
\min_{(\kappa,\phi)\in\mathcal{B}} J[\kappa,\phi].
\]
\end{definition}

\subsection{Interpretation}

\begin{remark}
This problem extends classical transition design in two directions:
\begin{itemize}
\item from pure geometry to coupled geometry--cant design,
\item from curve selection to functional optimization.
\end{itemize}
\end{remark}

\begin{remark}
In this framework, the transition curve is no longer chosen from a
catalogue, but emerges as part of an optimal coupled state.
\end{remark}

\subsection{Special Cases}

\begin{proposition}
If \(\phi(s)\equiv 0\), the coupled problem reduces to a pure curvature
optimization problem.
\end{proposition}

\begin{proposition}
If \(a_{\mathrm{eff}}(s)=0\) for all \(s\), then
\[
v^2 \kappa(s)=g\sin\phi(s),
\]
and the design is perfectly balanced.
\end{proposition}

\begin{remark}
These special cases show that classical equilibrium design appears as a
subcase of the coupled optimization framework.
\end{remark}

\subsection{Relation to Optimal Control}

\begin{remark}
The coupled optimization problem may be interpreted as an optimal
control problem with controls
\[
\kappa(s), \quad \phi'(s),
\]
state variables
\[
x(s),y(s),\theta(s),\phi(s),
\]
and cost functional \(J[\kappa,\phi]\).
\end{remark}

\begin{remark}
This provides a direct bridge from railway alignment design to modern
control-theoretic and variational methods.
\end{remark}

\subsection{Connection to the Main Thesis}

\begin{summary}
The coupled optimization problem shows that railway alignment design is
not a purely geometric task.

Instead, it is the optimization of a coupled state system in which
curvature and cant jointly determine dynamic behaviour.

This supports the central thesis of the manuscript that transition
curves and cant evolution should be understood as control functions in a
curvature-based railway state model.
\end{summary}

\section{Summary}

\begin{summary}
Alignment optimization treats geometry as a variable rather than a fixed
input.

It combines parameterization, numerical reconstruction, and constrained
optimization methods to produce alignments that satisfy both geometric
and engineering requirements.
\end{summary}

\section{Discrete Formulation of Alignment Optimization}

\subsection{Discretization of the Domain}

\begin{definition}[Discretization]
Let
\[
0 = s_0 < s_1 < \dots < s_N = L
\]
be a partition of the interval \([0,L]\).
\end{definition}

\begin{remark}
We denote
\[
\Delta s_i = s_{i+1} - s_i.
\]
\end{remark}

---

\subsection{Discrete Variables}

\begin{definition}[Discrete curvature]
Let
\[
\kappa_i \approx \kappa(s_i)
\]
denote the curvature values at the nodes.
\end{definition}

\begin{definition}[Discrete cant]
Let
\[
\phi_i \approx \phi(s_i)
\]
denote the cant values.
\end{definition}

\begin{definition}[Parameter vector]
The optimization variable is
\[
\mathbf{x}
=
(\kappa_0,\dots,\kappa_N,\phi_0,\dots,\phi_N) \in \R^{2(N+1)}.
\]
\end{definition}

---

\subsection{Discrete Pose Reconstruction}

\begin{remark}
The pose3 state is reconstructed by numerical integration.
\end{remark}

\begin{definition}[Discrete integration]
Given initial conditions
\[
\gamma_0,\ \theta_0,\ \phi_0,
\]
we compute iteratively:
\[
\theta_{i+1} = \theta_i + \kappa_i \Delta s_i,
\]
\[
\gamma_{i+1} = \gamma_i + \Delta s_i
(\cos\theta_i,\sin\theta_i),
\]
\[
\phi_{i+1} = \phi_i + \omega_i \Delta s_i.
\]
\end{definition}

\begin{remark}
Here, \(\omega_i\) may be expressed via differences of \(\phi_i\).
\end{remark}

---

\subsection{Finite Differences}

\begin{definition}[First derivative]
Approximate derivatives by
\[
\kappa'(s_i)
\approx
\frac{\kappa_{i+1}-\kappa_i}{\Delta s_i},
\quad
\phi'(s_i)
\approx
\frac{\phi_{i+1}-\phi_i}{\Delta s_i}.
\]
\end{definition}

\begin{definition}[Second derivative]
\[
\kappa''(s_i)
\approx
\frac{\kappa_{i+1}-2\kappa_i+\kappa_{i-1}}{(\Delta s_i)^2}.
\]
\end{definition}

---

\subsection{Discrete Objective Functions}

\subsubsection{Curvature Smoothness}

\[
J_\kappa(\mathbf{x})
=
\sum_{i=0}^{N-1}
\left(
\frac{\kappa_{i+1}-\kappa_i}{\Delta s_i}
\right)^2
\Delta s_i.
\]

---

\subsubsection{Cant Smoothness}

\[
J_\phi(\mathbf{x})
=
\sum_{i=0}^{N-1}
\left(
\frac{\phi_{i+1}-\phi_i}{\Delta s_i}
\right)^2
\Delta s_i.
\]

---

\subsubsection{Dynamic Objective}

\[
J_{\mathrm{dyn}}(\mathbf{x})
=
\sum_{i=0}^{N}
\bigl(
v^2 \kappa_i - g \sin\phi_i
\bigr)^2
\Delta s_i.
\]

---

\subsubsection{Schwerpunkt Objective}

\begin{remark}
Using finite differences, we approximate
\[
\|\gamma_c''(s_i)\|^2
\]
by second-order differences of \(\gamma_c\).
\end{remark}

\[
J_{\mathrm{veh}}(\mathbf{x})
=
\sum_{i=1}^{N-1}
\|\gamma_{c,i+1} - 2\gamma_{c,i} + \gamma_{c,i-1}\|^2.
\]

---

\subsection{Combined Discrete Problem}

\begin{definition}[Discrete objective]
\[
J(\mathbf{x})
=
\alpha J_{\mathrm{veh}}
+
\beta J_{\mathrm{dyn}}
+
\gamma J_\kappa
+
\delta J_\phi.
\]
\end{definition}

---

\subsection{Constraints}

\begin{definition}[Equality constraints]
\[
\kappa_0 = 0,
\quad
\kappa_N = \kappa_1.
\]
\end{definition}

\begin{definition}[Inequality constraints]
\[
|\kappa_i| \le \kappa_{\max},
\quad
|\phi_i| \le \phi_{\max}.
\]
\end{definition}

\begin{remark}
Further constraints may include:
\begin{itemize}
\item bounds on derivatives,
\item continuity conditions,
\item alignment anchoring points.
\end{itemize}
\end{remark}

---

\subsection{Final Optimization Problem}

\begin{definition}[Discrete optimization problem]
Find
\[
\mathbf{x}^\ast \in \R^{2(N+1)}
\]
such that
\[
J(\mathbf{x}^\ast)
=
\min_{\mathbf{x}} J(\mathbf{x})
\]
subject to the given constraints.
\end{definition}

---

\subsection{Structure of the Problem}

\begin{remark}
The problem has the structure of a nonlinear least-squares problem.
\end{remark}

\begin{remark}
It is therefore well suited for:
\begin{itemize}
\item Sequential Quadratic Programming,
\item Gauss--Newton methods,
\item sparse optimization techniques.
\end{itemize}
\end{remark}

---

\subsection{Connection to Implementation}

\begin{remark}
The discrete formulation directly corresponds to a computational model:
\begin{itemize}
\item \(\kappa_i\) correspond to element parameters,
\item \(\phi_i\) correspond to cant samples,
\item reconstruction matches numerical evaluation in software.
\end{itemize}
\end{remark}

\begin{summary}
The continuous alignment optimization problem can be reduced to a finite
dimensional nonlinear optimization problem.

This provides the direct link between theoretical modelling and
practical implementation.
\end{summary}

\section{Sequential Quadratic Programming (Concrete Formulation)}

\subsection{Problem Structure}

\begin{remark}
We consider the discrete optimization problem
\[
\min_{\mathbf{x}} J(\mathbf{x})
\quad \text{subject to} \quad
g(\mathbf{x}) = 0,\;
h(\mathbf{x}) \le 0.
\]
\end{remark}

\begin{remark}
The variable vector is
\[
\mathbf{x} =
(\kappa_0,\dots,\kappa_N,\phi_0,\dots,\phi_N).
\]
\end{remark}

---

\subsection{Residual Formulation}

\begin{definition}[Residual vector]
The objective function is written as a sum of squared residuals:
\[
J(\mathbf{x}) = \frac{1}{2} \| \mathbf{r}(\mathbf{x}) \|^2.
\]
\end{definition}

\begin{remark}
Typical residual components are:
\[
r_i^{(\kappa)} = \sqrt{\gamma}\;
\frac{\kappa_{i+1}-\kappa_i}{\Delta s},
\]
\[
r_i^{(\phi)} = \sqrt{\delta}\;
\frac{\phi_{i+1}-\phi_i}{\Delta s},
\]
\[
r_i^{(\mathrm{dyn})}
=
\sqrt{\beta}\;
(v^2 \kappa_i - g \sin\phi_i).
\]
\end{remark}

---

\subsection{Jacobian Matrix}

\begin{definition}[Jacobian]
The Jacobian of the residual is
\[
J_r(\mathbf{x})
=
\frac{\partial \mathbf{r}}{\partial \mathbf{x}}.
\]
\end{definition}

\begin{remark}
Each residual depends only on local variables, therefore the Jacobian is sparse.
\end{remark}

\begin{example}
For the curvature smoothness term:
\[
\frac{\partial r_i^{(\kappa)}}{\partial \kappa_i}
=
-\frac{\sqrt{\gamma}}{\Delta s},
\quad
\frac{\partial r_i^{(\kappa)}}{\partial \kappa_{i+1}}
=
\frac{\sqrt{\gamma}}{\Delta s}.
\]
\end{example}

---

\subsection{Quadratic Subproblem}

\begin{definition}[SQP step]
At iteration \(k\), we compute an increment \(\mathbf{p}\) by solving:
\[
\min_{\mathbf{p}}
\quad
\frac{1}{2} \mathbf{p}^T H_k \mathbf{p}
+
\nabla J(\mathbf{x}_k)^T \mathbf{p}
\]
subject to linearized constraints:
\[
g(\mathbf{x}_k) + \nabla g(\mathbf{x}_k)\mathbf{p} = 0.
\]
\end{definition}

---

\subsection{Gauss--Newton Approximation}

\begin{remark}
For least-squares problems, the Hessian is approximated by
\[
H_k \approx J_r^T J_r.
\]
\end{remark}

\begin{remark}
This leads to the Gauss--Newton system:
\[
(J_r^T J_r)\mathbf{p}
=
- J_r^T \mathbf{r}.
\]
\end{remark}

---

\subsection{Sparse Structure}

\begin{remark}
The matrix \(J_r^T J_r\) has band structure:
\begin{itemize}
\item curvature terms couple \(\kappa_i, \kappa_{i+1}\),
\item cant terms couple \(\phi_i, \phi_{i+1}\),
\item dynamic terms are diagonal.
\end{itemize}
\end{remark}

\begin{remark}
This enables efficient solution using sparse linear solvers.
\end{remark}

---

\subsection{Constraints Handling}

\begin{remark}
Constraints can be handled via:
\begin{itemize}
\item Lagrange multipliers,
\item penalty methods,
\item projection methods.
\end{itemize}
\end{remark}

\begin{example}
Simple bound constraints:
\[
\kappa_i \leftarrow \mathrm{clip}(\kappa_i, -\kappa_{\max}, \kappa_{\max}).
\]
\end{example}

---

\subsection{Iteration Scheme}

\begin{algorithm}
\caption{SQP iteration}
\begin{algorithmic}
\State initialize $\mathbf{x}_0$
\For{$k = 0,1,\dots$}
  \State compute residual $\mathbf{r}(\mathbf{x}_k)$
  \State compute Jacobian $J_r(\mathbf{x}_k)$
  \State solve
  \[
  (J_r^T J_r)\mathbf{p}_k = -J_r^T \mathbf{r}
  \]
  \State update
  \[
  \mathbf{x}_{k+1} = \mathbf{x}_k + \alpha_k \mathbf{p}_k
  \]
\EndFor
\end{algorithmic}
\end{algorithm}

---

\subsection{Step Size Control}

\begin{remark}
The step size \(\alpha_k\) is determined by:
\begin{itemize}
\item line search,
\item trust-region methods.
\end{itemize}
\end{remark}

---

\subsection{Connection to Implementation}

\begin{remark}
The SQP scheme maps directly to code:
\begin{itemize}
\item residual computation → evaluation functions,
\item Jacobian → local finite differences or analytic expressions,
\item linear solve → numeric library,
\item update → parameter array.
\end{itemize}
\end{remark}

---

\subsection{Interpretation}

\begin{summary}
Sequential Quadratic Programming provides a practical method for solving
the discrete alignment optimization problem.

Due to the sparse structure and local coupling, the problem is highly
efficiently solvable and directly implementable.

This establishes a concrete computational pathway from the theoretical
framework to a working alignment optimization engine.
\end{summary}

\section{Construction of Initial Guesses}

\subsection{Motivation}

\begin{remark}
The performance of SQP methods strongly depends on the choice of the
initial parameter vector.
\end{remark}

\begin{remark}
In railway applications, initial data is often available in the form of:
\begin{itemize}
\item measurement polylines,
\item existing alignments,
\item partially structured design data.
\end{itemize}
\end{remark}

---

\subsection{Input Representation}

\begin{definition}[Input polyline]
Let
\[
P = (p_0,\dots,p_M), \quad p_i \in \R^2
\]
denote a measured or imported polyline.
\end{definition}

---

\subsection{Arc-Length Parameterization}

\begin{definition}
Define cumulative arc-length
\[
s_0 = 0,
\quad
s_{i+1} = s_i + \|p_{i+1} - p_i\|.
\]
\end{definition}

\begin{remark}
This defines a parameterization of the input data.
\end{remark}

---

\subsection{Initial Curvature Estimation}

\begin{definition}[Discrete curvature]
For three consecutive points, define:
\[
\kappa_i
\approx
\frac{\theta_{i+1} - \theta_i}{\Delta s_i},
\]
where
\[
\theta_i = \arg(p_{i+1} - p_i).
\]
\end{definition}

\begin{remark}
This provides a first estimate of curvature from geometric data.
\end{remark}

---

\subsection{Smoothing}

\begin{remark}
Raw curvature estimates are typically noisy.
\end{remark}

\begin{definition}[Smoothing operator]
Apply a smoothing filter:
\[
\kappa_i^{(0)}
=
\sum_j w_{ij} \kappa_j.
\]
\end{definition}

\begin{remark}
Possible choices:
\begin{itemize}
\item moving average,
\item Gaussian filter,
\item low-pass filtering.
\end{itemize}
\end{remark}

---

\subsection{Initial Cant Estimation}

\begin{remark}
Cant may be initialized based on equilibrium:
\[
\phi_i^{(0)}
=
\arcsin\left(\frac{v^2 \kappa_i^{(0)}}{g}\right).
\]
\end{remark}

\begin{remark}
Alternatively, cant may be:
\begin{itemize}
\item set to zero,
\item taken from input data,
\item approximated by design rules.
\end{itemize}
\end{remark}

---

\subsection{Projection to Discrete Grid}

\begin{definition}
Interpolate the values \(\kappa_i^{(0)},\phi_i^{(0)}\)
onto the optimization grid:
\[
s_0,\dots,s_N.
\]
\end{definition}

---

\subsection{Heuristic Segmentation}

\begin{remark}
The input alignment may be segmented into:
\begin{itemize}
\item straight sections,
\item circular arcs,
\item transition regions.
\end{itemize}
\end{remark}

\begin{definition}[Segment detection]
A segment is classified based on curvature:
\[
\kappa_i \approx 0 \quad \Rightarrow \text{straight},
\]
\[
\kappa_i \approx \text{const} \quad \Rightarrow \text{arc}.
\]
\end{definition}

\begin{remark}
This structure can be used to improve the initial guess.
\end{remark}

---

\subsection{Final Initial Guess}

\begin{definition}[Initial parameter vector]
The initial guess is
\[
\mathbf{x}_0 =
(\kappa_0^{(0)},\dots,\kappa_N^{(0)},
\phi_0^{(0)},\dots,\phi_N^{(0)}).
\]
\end{definition}

---

\subsection{Interpretation}

\begin{remark}
The initial guess combines:
\begin{itemize}
\item measured geometry,
\item smoothing,
\item physical approximation.
\end{itemize}
\end{remark}

\begin{summary}
A robust initial guess transforms raw geometric input into a structured
parameter vector suitable for optimization.

This step is essential for bridging real-world data and numerical
methods.
\end{summary}
